\setcounter{footnote}{81}

\printedpage{104}
% head: A. Grothendieck : Dérivateurs, Ch. VI, Édition des Universités de Montpellier II et Paris VII
% foot: Version 9 novembre 2001

\msp{page 171}

\subsection*{New version of the theory of $\mathrm{HOT}_W$.}

I realized that the hypothesis $W \supseteq W_\infty$ (L3 quater, p. 1), which I had ended up putting on $W$, is artificial, and that the hypothesis L5 (p. 43), or even L6, L6 bis (p. 43-44), is unnecessary for a large part of the theory of $\mathrm{HOT}_W$. I have thus been led to review the theory (i.e. what I have done with it up to now) under the following hypotheses on $W$ (cf. XII):

\begin{itemize}[leftmargin=4em]
\item[W(1)] (\textit{Saturation} of $W$): as before (L1 page 1)
\item[W(2)] (localization): this is the old L4 bis (\footnote{i.e. L5 of page 2. The axiom is taken up again (in a similar and equivalent form) at the bottom of page 43.}).
\item[W(3)] (final object): If $X$ has a final object, $X$ is $W$-aspheric (old L3).
\end{itemize}

These three axioms are sufficient for everything I had done in Pursuing Stacks, and in a large part of what has been done here on HOT. Sometimes however one needs the old L 5, which I restate here (compare XII, theorem \texttt{[theorem number missing]}): % [sic] pout -> pour, misprint

\begin{itemize}[leftmargin=4em]
\item[W(4b)] Let $X$, $Y$ be $W$-fibered over $S$, $f : X \longrightarrow Y$ an $S$-morphism. If $f \in W$, then $\forall s \in \mathrm{Ob}\, S$, $f_s : X_s \longrightarrow Y_s$ is in $W$.
\end{itemize}

For the notion of $W$-fibered (\textit{weak}), I took a long time finally to find \textit{the} right notion, right for a $W$ satisfying only $\mathrm{W}(1,2,3)$ (and even only $\mathrm{W}(1, 2 \text{ bis}, 3)$, where W(2 bis) is the old L4). For this I refer to (XII, theorem 3, p. 51), where the notion is well pinned down by means of equivalent conditions. In fact, for the main applications to $\mathrm{HOT}_W$, seen up to now, one needs only (\footnote{This is ultimately an error, it seems now. I think that the axiom W(4b$'$) ultimately serves strictly no purpose (except in fact to keep false ideas alive.)}) a weakened variant of W(4b):

\begin{itemize}[leftmargin=4em]
\item[W(4b$'$)] Like W(4b), but requiring in addition that $X$, $Y$ be in $\ul{Parf}_\omega\, S$.
\end{itemize}

\msp{page 172}

For the notion of $W$-proper, $W$-smooth, $W$-perfect morphisms in Cat, I refer to XII, part III (p. 22-52), and the notations $\ul{Liss}_W S$, $\ul{Prop}_W S$, $\ul{Parf}_W S$ (subcategories of $\mathrm{Cat}/S$), and $\mathrm{Parf}_W \subseteq \mathrm{Fl}(\mathrm{Cat})$. In the statement of W(4b$'$) we have set
\[
\ul{Parf}_\omega S = \ul{Parf}_{W_\omega} S, \quad
\begin{array}[t]{l}
\text{where } W_\omega \subseteq \mathrm{Fl}(\mathrm{Cat}) \text{ is the smallest subset of } \mathrm{Fl}(\mathrm{Cat})\\
\text{satisfying } \mathrm{W}(1, 2\mathrm{bis}, 3).
\end{array}
\]

I similarly write
\[
W'_\omega \subseteq \mathrm{Fl}(\mathrm{Cat}), \qquad \text{smallest subset of } \mathrm{Fl}(\mathrm{Cat}) \text{ satisfying } \mathrm{W}(1, 2, 3, 4\mathrm{b}). % [sic] satifaisant -> satisfaisant, misprint
\]

\printedpage{105}
% head: A. Grothendieck : Dérivateurs, Ch. VI, Édition des Universités de Montpellier II et Paris VII
% foot: Version 9 novembre 2001

The hypothesis $W \supseteq W_\infty$ which we had ended up imposing on $W$, and which appeared artificial in the course of the reflection in XII, is better replaced by % [sic] apparu -> apparue, agreement with « l'hypothèse »
\[
W \supseteq W'_\omega \qquad (\text{or by } W \supseteq \overline{W'_\omega}, \text{ where } \overline{W'_\omega} \text{ is the strong saturation of } W'_\omega)
\]
(\footnote{$\overline{W'_\omega}$ is the smallest \textit{strongly saturated} subset of $\mathrm{Fl}(\mathrm{Cat})$, satisfying W(4b), and it even satisfies condition W$'$(4a), the strongest variant of the axioms of type W(4).}). But we do not impose it as a mandatory condition, on the same footing as $\mathrm{W}(1,2,3)$, any more than W(4b). The role I have had $W_\infty$ play in the past weeks, seems to belong more naturally to $W'_\omega$ or $\overline{W'_\omega}$. This is so in particular also, it seems to me, in the statement of the conditions of type Der 3 on a derivator (I, p. 111-112), which ought to be reviewed in this light. But everything depends, whether (as is my impression) one can really do substantial homotopical algebra, with in particular the long exact sequences of fibration and suspension, for $\mathrm{HOT}_W$, with $W$'s which are \textit{strictly} contained in $W_\infty$. Before even being sure, it is still too early for me to pronounce on the ``right'' formulation of Der 3,

\msp{page 173}

which is the axiom relating the derivator $\mathbf{D}$ to $W_\infty$, or to some other basic localizer.

I am now going to give a summary of the theory of $\mathrm{HOT}_W$, reviewed from the viewpoint I have just stated: we assume for $W$ \textit{the conditions} $\mathrm{W}(1,2,3)$ \textit{and no more,} unless explicitly mentioned.

\circled{1} \textbf{Main localizers} on $\mathrm{Cat}/S$ and on the main categories of $S$-categories (cf. diagram p. 76).

They will all be induced by the three familiar localizers
\[
\begin{tikzcd}[column sep=small]
W_S^{\mathrm{d}} \arrow[dr, hook] & & W_S^{\mathrm{g}} \arrow[dl, hook'] \\
& W_S &
\end{tikzcd}
\]
on $\mathrm{Cat}/S$. There is also reason to look at $W_S^{\mathrm{f}}$, but never on the whole of $\mathrm{Cat}/S$, but on other categories of the diagram, and always in cases where $W_S^{\mathrm{f}}$ coincides either with $W_S^{\mathrm{d}}$ or with $W_S^{\mathrm{g}}$. But the $W_S^{\mathrm{f}}$ point of view is often more directly convenient, and moreover has the advantage of being symmetric.

\circled{a} In $\ul{Liss}_W S$, we have

\printedpage{106}
% head: A. Grothendieck : Dérivateurs, Ch. VI, Édition des Universités de Montpellier II et Paris VII
% foot: Version 9 novembre 2001

\[
\begin{tikzcd}[row sep=large, column sep=0pt]
W_S^{\mathrm{f}} \arrow[d, Rightarrow] & {} = W_S^{\mathrm{d}} = W_S^{\mathrm{univ}} \\
W_S^{\mathrm{g}} \arrow[d, Rightarrow] & \\
W_S &
\end{tikzcd}
\]

Dually,

\circled{b} in $\ul{Prop}_W S$, we have
\[
\begin{tikzcd}[row sep=large, column sep=0pt]
W_S^{\mathrm{f}} \arrow[d, Rightarrow] & {} = W_S^{\mathrm{g}} = W_S^{\mathrm{univ}} \\
W_S^{\mathrm{d}} \arrow[d, Rightarrow] & \\
W_S, &
\end{tikzcd}
\]

\msp{page 174}

finally,

\circled{c} in $\ul{Fib}_W S$, we have
\[
\begin{tikzcd}[row sep=large]
W_S^{\mathrm{f}} \arrow[r, Leftrightarrow] \arrow[d, Rightarrow] & W_S^{\mathrm{d}} \arrow[r, Leftrightarrow] & W_S^{\mathrm{g}} \arrow[r, Leftrightarrow] & W_S^{\mathrm{univ}} \\
W_S. & & &
\end{tikzcd}
\]

This implication $W_S^{\mathrm{f}} \Longrightarrow W_S$ is strict if one does not impose on $W$ axiom W(4b), whose usefulness here is exactly to ensure
\[
\left|\;
\begin{array}{l}
W_S^{\mathrm{f}} \Longleftrightarrow W_S \ \text{ in } \ul{Fib}_W S \ \textit{if}\ \mathrm{W}(4\mathrm{b})\ \texttt{[is]}\ \text{satisfied}\\[1.5ex]
\text{and then the five localizers considered above}\\[1.5ex]
\text{in } \ul{Fib}_W S \text{ are all equal.}
\end{array}
\right.
\]

The relations between [?] localizers in $\ul{Liss}_W S$, $\ul{Prop}_W S$, $\ul{Fib}_W S$ carry over automatically to subcategories of these, more generally to all those which map
\printedpage{107}
% head: A. Grothendieck : Dérivateurs, Ch. VI, Édition des Universités de Montpellier II et Paris VII
% foot: Version 9 novembre 2001
by an inclusion functor into one of these. In the diagram p. 76, all the categories, with the sole exception of the whole of $\mathrm{Cat}/S$, are in this case.

For the main results developed up to now on $\mathrm{HOT}_W$, when axiom W(4b) is used, the weak form W(4b$'$) suffices. We have
\[
\left|\;
\begin{array}{l}
W_S^{\mathrm{f}} \Longleftrightarrow W_S \ \text{ in } \ul{Parf}_\omega S \ \textit{if}\ \mathrm{W}(4\mathrm{b}')\ \texttt{[is]}\ \text{satisfied}\\[1.5ex]
\text{and then the five localizers } W_S^{\mathrm{f}},\, W_S^{\mathrm{d}},\, W_S^{\mathrm{g}},\, W_S^{\mathrm{univ}},\, W_S\\[1.5ex]
\text{are equal in } \ul{Parf}_\omega S.
\end{array}
\right.
\]

\circled{2} \textbf{Main categories of categories over $S$.}

These are, up to changes in lighting, those of the diagram p. 76. The categories $\ul{Fibsc}\,S$, $\ul{Fib}\,S$, $\ul{Cofibsc}\,S$, $\ul{Cofib}\,S$ and their variants with subscript 0, are unchanged - their role is to be as small

\msp{page 175}

as possible to fulfill their function (of describing $\mathrm{HOT}_W(S)$, $\mathrm{HOT}^{o}_W(S)$, $\mathrm{HOT}^{\mathrm{lc}}_W(S)$, or $\mathrm{HOT}^{\mathrm{lc}!}_W(S)$ below). The category $\ul{Parf}\,S$ of the diagram fulfilled a function of the same type, which is why we now replace it (bidding higher) by $\ul{Parf}_\omega S$. The other categories of the diagram are, in addition to $\mathrm{Cat}/S$, $\ul{Liss}\,S$ and $\ul{Prop}\,S$ and their variants with subscript 0, as well as $\ul{Fibhot}\,S$. (All these categories, as well as $\ul{Parf}\,S$, were thus constructed with reference to $W_\infty$, not to $W$.) Their role is to be as large as possible to give rise to certain properties and relations. Now it turns out that here the reference to $W_\infty$ was artificial, and one should put instead the categories $\ul{Liss}_W S$, $\ul{Prop}_W S$, $\ul{Fib}_W S$, and in addition the variants $\ul{Liss}_{0,W} S$, $\ul{Prop}_{0,W} S$. Thus the diagram becomes % [sic] propiétés -> propriétés, misprint
\[
\begin{tikzcd}[column sep=tiny, row sep=small]
 & & & & \ul{Fib}_W S \arrow[dd, hook, "\delta"] \arrow[dddllll, dotted, bend right=35, "\Phi_0"'] \arrow[dddrrrr, dotted, bend left=35, "\Psi_0"] & & & & \\
 & & \ul{Liss}_{0,W} S \arrow[urr, hook] \arrow[dd, hook] & & & & \ul{Prop}_{0,W} S \arrow[ull, hook'] \arrow[dd, hook] & & \\
 & \ul{Fib}_{0,W} S \arrow[ur, hook] \arrow[dd, hook] & & & \mathrm{Cat}/S \arrow[dd, "C"] \arrow[dddllll, dotted, bend right=30, "\Phi"] \arrow[dddrrrr, dotted, bend left=30, "\Psi"'] & & & \ul{Cofib}_{0,W} S \arrow[ul, hook'] \arrow[dd, hook] & \\
\ul{Fibsc}_{0,W} S \arrow[ur] \arrow[dd, hook] & & \ul{Liss}_W S \arrow[urr, hook] & & & & \ul{Prop}_W S \arrow[ull, hook'] & & \ul{Cofibsc}_{0,W} S \arrow[ul] \arrow[dd, hook] \\
 & \ul{Fib}\,S \arrow[ur, hook] & & & \ul{Parf}_\omega S \arrow[uuull, hook] \arrow[uuurr, hook'] & & & \ul{Cofib}\,S \arrow[ul, hook'] & \\
\ul{Fibsc}\,S \arrow[ur] & & & & & & & & \ul{Cofibsc}\,S \arrow[ul]
\end{tikzcd}
\]
(\footnote{NB the new notation would rather be $\ul{Fibf}_W S$, for these are $W$-fibrations which I now call ``weak'' (and defined only in Cat, as opposed to $W$-fibrations, definable in any category of models).}) (all the squares \texttt{[are]} cartesian, the functors denoted $\hookrightarrow$ are fully faithful),

\printedpage{108}
% head: A. Grothendieck : Dérivateurs, Ch. VI, Édition des Universités de Montpellier II et Paris VII
% foot: Version 9 novembre 2001

\msp{page 176}

where I have again marked in green \texttt{[in the original only]} the five \textit{``inversion functors''}
\[
\Phi,\ \Psi,\ \Phi_0,\ \Psi_0,\ C,
\]
which provide quasi-inverses of the functors on the localized categories induced by the \textit{inclusion functors.}

\circled{3} \textbf{The four types of categories} $\mathrm{HOT}_W$\textbf{.}

On each of the categories of the upper floor, there are \textit{two} localizers to look at, % [sic] l'etage -> l'étage, missing accent (printed « l'étage » below)
\[
\begin{tikzcd}[row sep=huge, column sep=0pt]
W_S^{\mathrm{f}} \arrow[d, Rightarrow] & {} = W_S^{\mathrm{g}} = W_S^{\mathrm{d}} = W_S^{\mathrm{univ}} \\
W_S, &
\end{tikzcd}
\]
giving rise to two distinct localized categories, denoted
\[
\mathrm{HOT}^{\mathrm{lc}}_W(S) \qquad \text{and} \qquad \mathrm{HOT}^{\mathrm{lc}!}_W(S),
\]
it turns out that (thanks to $\Phi_0$ and $\Psi_0$) these localized categories ``are the same'' for the 8 categories of the upper floor, with the exception of $\ul{Parf}_\omega S$; more exactly, the inclusion functors (with the exception perhaps of those having as source $\ul{Parf}_\omega S$) induce equivalences on the two types of localized categories, the quasi-inverse equivalences being obtained by $\Phi_0$ (left wing) or $\Psi_0$ (right wing). When one localizes \textit{for} $W_S$ (to the exclusion of $W_S^{\mathrm{f}}$), one also finds quasi-inverses by $C|\ul{Liss}_{0,W} S$ resp. $C|\ul{Prop}_{0,W} S$, for the two inclusions
\[
\begin{tikzcd}
\ul{Liss}_{0,W} S & \ul{Parf}_\omega S \arrow[l, hook'] \arrow[r, hook] & \ul{Prop}_{0,W} S
\end{tikzcd}
\]
(\footnote{cf. below, what one can say for the localized categories for $W_S^{\mathrm{f}}$.}). For reasons of symmetry, the prettiest seems to be to define
\[
\left\{
\begin{array}{lcl}
\mathrm{HOT}^{\mathrm{lc}}_W(S) & = & (\ul{Fib}_W S)(W_S^{\mathrm{f}})^{-1}\\
\mathrm{HOT}^{\mathrm{lc}!}_W(S) & = & (\ul{Fib}_W S)(W_S)^{-1},
\end{array}
\right.
\]
and we have a localization functor
\[
\mathrm{HOT}^{\mathrm{lc}}_W(S) \xrightarrow{\ \lambda\ } \mathrm{HOT}^{\mathrm{lc}!}_W(S)
\]
induced by the inclusion
\[
W_S^{\mathrm{f}} \subseteq W_S.
\]